\section{Fiche de suivi et d'évaluation N1}
\label{sec:fiche-evaluation}

Cette fiche reprend la proposition de fiche de suivi et d'évaluation N1 du manuel de
formation FFESSM (PE20), en détaillant chaque compétence par ses \textbf{sous-compétences}
techniques afin de permettre un suivi plus précis, séance après séance (colonnes S1 à S8).

\textbf{Nom :} \makebox[5cm]{\hrulefill} \hfill \textbf{Prénom :} \makebox[5cm]{\hrulefill}

\begin{landscape}
  \footnotesize
  \renewcommand{\arraystretch}{1.25}
  \begin{longtable}{|p{3.2cm}|p{5.3cm}|c|c|c|c|c|c|c|c|}

    \caption{Fiche de suivi et d'évaluation détaillée --- Plongeur Niveau 1 (PE20)}\\
    \hline
    \rowcolor{ffessmblue}
    \textcolor{white}{\textbf{Compétence}} & \textcolor{white}{\textbf{Sous-compétence}} &
    \textcolor{white}{\textbf{S1}} & \textcolor{white}{\textbf{S2}} &
    \textcolor{white}{\textbf{S3}} & \textcolor{white}{\textbf{S4}} &
    \textcolor{white}{\textbf{S5}} & \textcolor{white}{\textbf{S6}} &
    \textcolor{white}{\textbf{S7}} & \textcolor{white}{\textbf{S8}} \\
    \hline
    \endfirsthead

    \multicolumn{10}{l}{\small\itshape ... suite de la page précédente}\\
    \hline
    \rowcolor{ffessmblue}
    \textcolor{white}{\textbf{Compétence}} & \textcolor{white}{\textbf{Sous-compétence}} &
    \textcolor{white}{\textbf{S1}} & \textcolor{white}{\textbf{S2}} &
    \textcolor{white}{\textbf{S3}} & \textcolor{white}{\textbf{S4}} &
    \textcolor{white}{\textbf{S5}} & \textcolor{white}{\textbf{S6}} &
    \textcolor{white}{\textbf{S7}} & \textcolor{white}{\textbf{S8}} \\
    \hline
    \endhead

    \hline
    \multicolumn{10}{r}{\small\itshape suite page suivante ...}\\
    \endfoot

    \hline
    \endlastfoot

    % --- S'équiper et se déséquiper ---
    \multirow{3}{=}{S'équiper et se déséquiper} & Gréage et dégréage & & & & & & & & \\ \cline{2-10}
    & Capelage et décapelage & & & & & & & & \\ \cline{2-10}
    & Choix de son matériel personnel & & & & & & & & \\
    \hline

    % --- Se mettre à l'eau et sortir de l'eau ---
    \multirow{4}{=}{Se mettre à l'eau et sortir de l'eau} & Saut droit & & & & & & & & \\ \cline{2-10}
    & Bascule arrière & & & & & & & & \\ \cline{2-10}
    & Départ plage & & & & & & & & \\ \cline{2-10}
    & Sortir de l'eau & & & & & & & & \\
    \hline

    % --- S'immerger ---
    \multirow{2}{=}{Évoluer dans l'eau --- S'immerger} & Canard & & & & & & & & \\ \cline{2-10}
    & Phoque & & & & & & & & \\
    \hline

    % --- Se propulser ---
    \multirow{5}{=}{Évoluer dans l'eau --- Se propulser} & Palmage ventral en surface & & & & & & & & \\ \cline{2-10}
    & Palmage dorsal & & & & & & & & \\ \cline{2-10}
    & Palmage de sustentation & & & & & & & & \\ \cline{2-10}
    & Palmage en immersion & & & & & & & & \\ \cline{2-10}
    & Nage en capelé & & & & & & & & \\
    \hline

    % --- Se ventiler ---
    \multirow{4}{=}{Évoluer dans l'eau --- Se ventiler} & Ventilation en immersion & & & & & & & & \\ \cline{2-10}
    & Ventilation sur tuba et vidage du tuba & & & & & & & & \\ \cline{2-10}
    & Vidage du masque & & & & & & & & \\ \cline{2-10}
    & Lâcher et reprise d'embout & & & & & & & & \\
    \hline

    % --- S'équilibrer ---
    \multirow{2}{=}{Évoluer dans l'eau --- S'équilibrer} & Gestion du gilet de stabilisation & & & & & & & & \\ \cline{2-10}
    & Poumon ballast & & & & & & & & \\
    \hline

    % --- Respecter le milieu et l'environnement ---
    Respecter le milieu et l'environnement & Aisance aquatique & & & & & & & & \\
    \hline

    % --- Communiquer ---
    Communiquer & Exécution des signes conventionnels & & & & & & & & \\
    \hline

    % --- Retourner en surface ---
    \multirow{5}{=}{Retourner en surface} & Maîtrise de la vitesse de remontée & & & & & & & & \\ \cline{2-10}
    & Tenue d'un palier & & & & & & & & \\ \cline{2-10}
    & Tour d'horizon & & & & & & & & \\ \cline{2-10}
    & Gonflage du gilet en surface & & & & & & & & \\ \cline{2-10}
    & Remontée en expiration contrôlée & & & & & & & & \\
    \hline

    % --- Évoluer en sécurité ---
    \multirow{2}{=}{Évoluer en sécurité} & Application des procédures mises en œuvre par le GP & & & & & & & & \\ \cline{2-10}
    & Intervention en relais & & & & & & & & \\
    \hline

  \end{longtable}
\end{landscape}

L'évaluation des compétences se fait en \textbf{contrôle continu}. Le niveau est validé
lorsque toutes les compétences sont acquises. Les connaissances théoriques sont évaluées à
l'oral, lors des mises en situations pratiques. Il n'y a pas d'examen écrit. L'accent est mis
sur la prévention.

\vspace{0.5cm}
\textbf{Validé le :} \makebox[4cm]{\hrulefill} \hfill
\textbf{Nom et signature du moniteur :} \makebox[6cm]{\hrulefill}

\vspace{0.3cm}
{\small\itshape Cette fiche n'a aucun caractère obligatoire, elle est modulable et il
appartient au moniteur de l'adapter pour sa propre utilisation.}
